\chapter{Caching}

Caching is the main lever for speed.
It also introduces correctness risks if you cache the wrong things.

\section{What to cache}
Cache build inputs, not outputs.
Examples:
\begin{itemize}[nosep]
  \item dependency downloads (\code{pip}, \code{poetry}, \code{yarn})
  \item build tool caches (Gradle, Cargo)
\end{itemize}

\section{Cache with setup actions}
Many setup actions support caching directly.

\begin{lstlisting}[language=YAML]
- uses: actions/setup-node@v4
  with:
    node-version: "20"
    cache: "yarn"
\end{lstlisting}

\section{Cache keys}
Keys should change when dependencies change.
The simplest reliable approach is hashing lockfiles.

\begin{lstlisting}[language=YAML]
- uses: actions/cache@v4
  with:
    path: |
      ~/.cache/pip
    key: pip-${{ runner.os }}-${{ hashFiles('**/requirements.txt') }}
\end{lstlisting}

\begin{warningbox}{Warning: caching virtual environments}
Caching whole virtualenv folders often causes subtle corruption across runner images.
Prefer caching download directories and rebuilding environments.
\end{warningbox}

\begin{checklistbox}{Checklist: safe caching}
\begin{itemize}[nosep]
  \item Cache by lockfile hash.
  \item Keep caches OS-specific.
  \item Prefer restoring caches as an optimization, not a requirement.
\end{itemize}
\end{checklistbox}
